% !TeX TS-program = pdflatex
% !BIB TS-program = biber
\documentclass[aspectratio=169,professionalfonts]{beamer}
\usepackage[libertine]{newtxmath}
\usepackage[osf]{ETbb}
\usepackage[colorlinks]{hyperref}
\usepackage{authblk,xspace}
\usepackage[margin=3cm,top=2cm]{geometry}
\setlength{\parskip}{1ex}
\setlength{\parindent}{0pt}
\newcommand{\hreftt}[1]{\href{#1}{{\footnotesize\texttt{#1}}}\xspace}
\def\covid{\textsc{covid-19}\xspace}
\def\dowc{\immediate\write18{detex abstract.tex | wc -w > wordcount.tex}}
\def\printwc{\input{wordcount.tex}}
\title{\vskip-8ex%
  Estimating effective reproduction number using
  generation time versus serial interval,\\
  with application to \covid\\
  in the Greater Toronto Area, Canada
}
\usepackage{authblk}
\author[1,2]{Jesse Knight}
\author[1,2,3,4]{Sharmistha Mishra}
\affil[1]{Institute of Medical Science, University of Toronto}
\affil[2]{MAP Centre for Urban Health Solutions, Unity Health Toronto}
\affil[3]{Division of Infectious Disease, Department of Medicine, University of Toronto}
\affil[4]{Institute of Health Policy, Management and Evaluation, Dalla Lana School of Public Health, University of Toronto}
\renewcommand\Affilfont{\footnotesize}
\date{
  {\textbf{Correction}}\\
  March 23, 2021\\[1ex]
  \parbox{\linewidth}{\footnotesizeIn a previous version of this work,
the parameters $\theta = [\alpha, \beta]$ were calculated incorrectly
because the Kullback-Leibler divergence was defined in the wrong direction.
The impact on generation time parameters and statistics is as follows
(original $\rightarrow$ fixed):
shape ($\alpha$): $1.813 \rightarrow 1.633$,
scale ($\beta$): $2.199 \rightarrow 2.498$,
mean: $3.99 \rightarrow 4.08$,
\sd: $2.96 \rightarrow 3.19$.
The qualitative interpretation of results is unchanged,
and the corrected numerical results are given in this version.
We sincerely apologize for this error.\\
The error correction is shown here:
\scriptsize\hreftt{https://github.com/mishra-lab/covid-r/commit/aaef512}
}
}

\usepackage[libertine]{newtxmath}
\usepackage[osf]{ETbb}
\usepackage[colorlinks]{hyperref}
\usepackage{authblk,xspace}
\usepackage[margin=3cm,top=2cm]{geometry}
\setlength{\parskip}{1ex}
\setlength{\parindent}{0pt}
\newcommand{\hreftt}[1]{\href{#1}{{\footnotesize\texttt{#1}}}\xspace}
\def\covid{\textsc{covid-19}\xspace}
\def\dowc{\immediate\write18{detex abstract.tex | wc -w > wordcount.tex}}
\def\printwc{\input{wordcount.tex}}
\begin{document}
% ------------------------------------------------------------------------------
\begin{frame}
  \maketitle
\end{frame}
% ------------------------------------------------------------------------------
\begin{frame}{Outline}
  \begin{itemize}
    \item Reproduction Number $R$
    \item Estimating $R$
    \item Estimating Generation Time Distribution $G(\tau)$
    \item Results
    \item Implications
  \end{itemize}
\end{frame}
% ==============================================================================
\section{Background}
% ------------------------------------------------------------------------------
\begin{frame}{Reproduction Number: Defined}
  Average number of secondary infections
  generated by each infectious individual
  \par
  \pause
  \centering\begin{tikzpicture}
\foreach \R in {1,2,3}{\visible<+>{
    \outbreak{2}{\R}{0}{0}{0}
    \node at ( 3,0){\dots};
    \node at (-1,0){$R = \R$};
}}
\end{tikzpicture}
\end{frame}
% ------------------------------------------------------------------------------
\begin{frame}{Basic $R_0$ and Effective $R_e(t)$}
  \begin{minipage}{0.6\linewidth}
    \begin{align*}
      \visible<1->{
        R_0    &= \beta \cdot C \cdot D}\\
      \visible<5->{
        R_e(t) &= \beta_t \cdot C_t \cdot D_t \cdot \hat{S}(t)}
    \end{align*}
    \begin{itemize}
      \item<2-> $\beta$ = probability of transmission per contact
      \item<3-> $C$ = contact rate
      \item<4-> $D$ = duration of infectiousness
      \item<6-> $\hat{S}$ = proportion who are susceptible
    \end{itemize}
  \end{minipage}%
  \begin{minipage}{0.4\linewidth}
    \visible<7>{\centering\begin{tikzpicture}
  \outbreak{2}{3}{1}{0}{0}
  \node at (-1.5,-1em){
    $\begin{aligned}
      R_e(t) &= 3 \cdot \frac{2}{3}\\ &= 2
    \end{aligned}$};
\end{tikzpicture}}
  \end{minipage}
\end{frame}
% ------------------------------------------------------------------------------
\begin{frame}{$R_0$ for Common Diseases}
  \centering
  \begin{tabular}{lcl}
	\toprule
	Disease      &      $R_0$       & Transmission \\
	\midrule
	Measles      &     12 -- 18     & Aerosol      \\
	Chickenpox   &     10 -- 12     & Aerosol      \\
	Polio        &      5 -- 7      & Fecal-Oral   \\
	Smallpox     &      3 -- 6      & Droplet      \\
	HIV          &      1 -- 5      & Body Fluids  \\
	Seasonal Flu &    1.5 -- 2.5    & Droplet      \\
	\midrule
	\clr{\covid} & \clr{1.9 -- 6.5}\footnotemark & Droplet \& Aerosol \\
	\bottomrule
\end{tabular}
\footnotetext{\cite{Park2020}}\quad
  \parbox{0.4\linewidth}{
    \begin{itemize}
      \item Without interventions
      \item Context-specific: $\beta \cdot C \cdot D$
  \end{itemize}}
\end{frame}
% ------------------------------------------------------------------------------
\begin{frame}{Applications of $R$}
  \begin{itemize}
    \pause
    \item Measure of epidemic ``potential''
    \pause
    \item Monitor interventions
    \pause
    \item Required vaccine coverage:
    \begin{equation*}
      V = 1 - \frac{1}{R_0}
    \end{equation*}
  \end{itemize}
\end{frame}
% ==============================================================================
\section{Methods}
% ------------------------------------------------------------------------------
\begin{frame}{Estimating $R$:\enspace3 Methods}
  \begin{itemize}
    \item Static exponential: $R_0$
    \item Renewal equation: $R_e(t)$
    \item Mechanistic models (not covered)
  \end{itemize}
\end{frame}
% ------------------------------------------------------------------------------
\begin{frame}{Exponential Model for $R_0$}
  \begin{minipage}{0.5\linewidth}
    \begin{align*}
      I(t) &= I_0 \cdot e^{\alpha t} \\
      \visible<3->{
        R_0  &= \frac{I(t+g)}{I(t)} = e^{\alpha g}}
    \end{align*}
    \begin{itemize}
      \item $I$ = infections
      \item $\alpha$ = exponential growth rate
      \item<2-> $g$ = generation time:\\
      \hspace{1.7em}time between infections
    \end{itemize}
  \end{minipage}%
  \begin{minipage}{0.5\linewidth}
    \centering
    \begin{tikzpicture}
  \pgfmathsetseed{123}
  \plot{C3}{\t}{0:5}{\t}{\texp{\kt}{\t}}{};
  \draw[thick,arrow,->] (0,0) -- node[anchor=north]{Time}                 (5.5,0);
  \draw[thick,arrow,->] (0,0) -- node[anchor=south,rotate=90]{Infections} (0,5);
  \foreach \t in {0,0.2,...,5}{
    \pgfmathsetmacro{\y}{max(0, random() - 0.5 + \texp{\kt}{\t})}
    \node[dot=C3] at (\t,\y){};
  }
  \pgfmathsetmacro{\yo}{\texp{\kt}{2.5}}
  \pgfmathsetmacro{\yf}{\texp{\kt}{4.0}}
  \draw[brace] (2.5,\yo) -- (4.0,\yo) node[midway,yshift=-0.4cm]{$g$}; % \visible<2->{
  \draw[brace] (4.0,\yo) -- (4.0,\yf) node[midway,xshift=+0.4cm]{$R_0$}; % \visible<3->{
\end{tikzpicture}

  \end{minipage}
\end{frame}
% ------------------------------------------------------------------------------
\begin{frame}{Exponential Model: Assumptions}
  \begin{minipage}{0.5\linewidth}
    \heading{Assumptions:} % TODO
    \begin{itemize}
      \item Homogeneous susceptible population
      \item Constant $R_0$
      \item Reported cases $\,\propto\,$ true cases
    \end{itemize}
    \vspace{1em}
    \heading{If longer generation time \textit{g}:}
    \begin{itemize}
      \item Expect slower epidemic growth
      \item Infer \emph{greater} $R_0$
    \end{itemize}
  \end{minipage}%
  \begin{minipage}{0.5\linewidth}
    \centering
    \begin{tikzpicture}
  \pgfmathsetseed{123}
  \plot{C3}{\t}{0:5}{\t}{\texp{\kt}{\t}}{};
  \draw[thick,arrow,->] (0,0) -- node[anchor=north]{Time}                 (5.5,0);
  \draw[thick,arrow,->] (0,0) -- node[anchor=south,rotate=90]{Infections} (0,5);
  \foreach \t in {0,0.2,...,5}{
    \pgfmathsetmacro{\y}{max(0, random() - 0.5 + \texp{\kt}{\t})}
    \node[dot=C3] at (\t,\y){};
  }
  \pgfmathsetmacro{\yo}{\texp{\kt}{2.5}}
  \pgfmathsetmacro{\yf}{\texp{\kt}{4.0}}
  \draw[brace] (2.5,\yo) -- (4.0,\yo) node[midway,yshift=-0.4cm]{$g$}; % \visible<2->{
  \draw[brace] (4.0,\yo) -- (4.0,\yf) node[midway,xshift=+0.4cm]{$R_0$}; % \visible<3->{
\end{tikzpicture}

  \end{minipage}
\end{frame}
% ------------------------------------------------------------------------------
\begin{frame}{Renewal Equation for $R_e(t)$}
  \invisible{\footcite{Fraser2007,Cori2013}}
  Similar to exponential model, but with generation time \emph{distribution}: $g \sim G(\tau)$
  \pause
  \begin{align*}
    I(t) &= R_e(t) \int_{t_0}^t I(t-\tau) G(\tau) d\tau\\
    \visible<3->{
      R_e(t) &= I(t) {\left[\int_{t_0}^{t}{I(t-\tau) G(\tau) d\tau}\right]}^{-1}}
  \end{align*}
  \begin{itemize}
    \item $I(t-h)$ = infections $\tau$ days ago
    \item $G(\tau)$ = infectiousness $\tau$ days since infection
    \item $t_0 = t-w$, integration window
  \end{itemize}
\end{frame}
% ------------------------------------------------------------------------------
\begin{frame}{Renewal Equation: Illustrated}
  \begin{minipage}{0.6\linewidth}
    \centering\begin{tikzpicture}
  \visible<2->{
    \def\ymin{3}
    \plot{C1}{\t}{0:5}{\t}{3+.2*(0
      +\erf{4.4}{.2}{\t}
      +\erf{3.2}{.2}{\t}
      +\erf{2.9}{.2}{\t}
      +\erf{2.8}{.2}{\t}
      +\erf{2.7}{.2}{\t}
      +\erf{2.3}{.2}{\t}
      +\erf{0.9}{.2}{\t}
      +\erf{0.6}{.2}{\t})
    }{};
    \draw[thick,arrow,->] (0,3) -- node[anchor=east] {\ulab{$\int_{t_0}^t I(t-\tau) G(\tau) d\tau$}} (0,5);
  }
  \foreach \i/\x in {8/4.4, 7/3.2, 6/2.9, 5/2.8, 4/2.7, 3/2.3, 2/0.9, 1/0.6}{
    \def\ymin{0.3*\i}
    \plot{C3}{\t}{0:5}{\t}{\ymin+.2*\gauss{\x}{.2}{\t}}{};
  }
  \draw[thick,arrow,->] (0,0) -- node[anchor=north]{time} (5.5,0);
  \draw[thick,brace] (-.3,2.7) -- node[left=.2\dx] {\ulab{$G_i(\tau)$}} (-.3,0.3);
\end{tikzpicture}
  \end{minipage}%
  \begin{minipage}{0.4\linewidth}
    \pause\pause
    \heading{Assumptions:}
    \begin{itemize}
      \item Same as exponential
      \item Constant $R_e$ over each window
    \end{itemize}
    \pause
    \vspace{1em}
    \heading{Moving window $\mathbf{\rightarrow  R_e(t)}$}
  \end{minipage}
\end{frame}
% ------------------------------------------------------------------------------
\begin{frame}{$R_e(t)$: Examples}
  \includegraphics[width=\linewidth]{Cori2013}
  \invisible{\footcite{Cori2013}}
\end{frame}
% ------------------------------------------------------------------------------
\begin{frame}{Infectiousness: Serial Interval $S(\tau)$}
  \begin{minipage}[t]{0.5\linewidth}
    \centering
    1. Known case-pairs:
    \\[1em]
    {\tiny
\begin{tabular}{ccc}
	\toprule
	         & \multicolumn{2}{c}{Symptom Onset} \\
	\cmidrule(rl){2-3}
	  Pair   &  Infector  &       Infectee       \\
	\midrule
	   1     & 2020-01-01 &      2020-01-07      \\
	   2     & 2020-01-04 &      2020-01-14      \\
	   3     & 2020-01-05 &      2020-01-09      \\
	   4     & 2020-01-02 &      2020-01-15      \\
	   5     & 2020-01-08 &      2020-01-12      \\
	   6     & 2020-01-12 &      2020-01-23      \\
	$\vdots$ &  $\vdots$  &       $\vdots$       \\
	\bottomrule
\end{tabular}
}
  \end{minipage}\pause%
  \begin{minipage}[t]{0.5\linewidth}
    \centering
    2. Fit a distribution:
    \\[1em]
    \begin{tikzpicture}
  \plot{C3}{\t}{0:5}{\t}{3*\gauss{2.3}{0.8}{\t}}{}
  \foreach \t in {0,0.25,...,5}{
    \pgfmathsetmacro{\y}{max(0,random() - 0.5 + 3*\gauss{2.3}{0.8}{\t})}
    \draw[draw=C2,fill=none] (\t,0) -- ++(0.25,0) -- ++(0,\y) -- ++(-0.25,0) -- cycle;
  }
  \draw[thick,arrow,->] (0,0) -- node[anchor=north]{Serial Interval}       (5.5,0);
  \draw[thick,arrow,->] (0,0) -- node[anchor=south,rotate=90]{Probability} (0,3);
\end{tikzpicture}
  \end{minipage}
\end{frame}
% ------------------------------------------------------------------------------
\begin{frame}{Infectiousness: Serial Interval vs Generation Time}
  \vspace{1em}
  \centering\only<2>{Serial Interval}
\only<3>{Generation Time}
\only<4>{Incubation Period $i$}
\only<5>{Incubation Period $i+1$}
\only<6>{\vphantom{I}}
\par
\scalebox{0.85}{
\begin{tikzpicture}
%  \node at (1,5.3){\smash{
%    \only<2>{Serial Interval}%
%    \only<3>{Generation Time}%
%    \only<4>{Incubation Period $i$}%
%    \only<5>{Incubation Period $i+1$}%
%    \only<6>{``undo'' incubation blur to recover $G(t)$ from $S(t)$}%
%  }};
  % nodes
  \visible<3,4,6>{\node (f1) [event=fr]       at ( 0\dx, 4\dx) {$f_i$};}
  \visible<3,5,6>{\node (f2) [event=to]       at ( 2\dx, 3\dx) {$f_{i+1}$};}
  \visible<2,4,6>{\node (y1) [event=fr!50!sy] at ( 0\dx, 2\dx) {$y_i$};}
  \visible<2,5,6>{\node (y2) [event=to!50!sy] at ( 2\dx, 1\dx) {$y_{i+1}$};}
  % arrows
  \visible<3,6>{\draw[arrow,->]        (f1) -- (f2);}
  \visible<2,6>{\draw[arrow,->,dashed] (y1) -- (y2);}
  \visible<4,6>{\draw[arrow,->]        (f1) -- (y1);}
  \visible<5,6>{\draw[arrow,->]        (f2) -- (y2);}
  \visible<  6>{\draw[arrow,->,dotted] (f2) -- ++( 2\dx,-1\dx);}
  \visible<  6>{\draw[arrow,<-,dotted] (f1) -- ++(-2\dx, 1\dx);}
  % distributions
  \visible<2->{
    \node(t0) at (5\dx, 5.3\dx) {\rlap{$t$}};
    \node(t1) at (5\dx,-0.5\dx) {};
  }
  \visible<3,6>{\plot{to}            {\t}{0:5}{{5\dx+0.5\dx*\gauss{3}{.23}{\t}}}{\t}{}}
  \visible<4,6>{\plot{fr!50!sy}      {\t}{0:5}{{5\dx+0.5\dx*\gauss{2}{.28}{\t}}}{\t}{}}
  \visible<2,6>{\plot{fr!50!to!50!sy}{\t}{0:5}{{5\dx+0.5\dx*\gauss{1}{.36}{\t}}}{\t}{}}
  \visible<5  >{\plot{to!50!sy}      {\t}{0:5}{{5\dx+0.5\dx*\gauss{1}{.23}{\t}}}{\t}{}}
  \visible<2->{\draw[thick,arrow,->] (t0) -- (t1);}
  \visible<2    >{\draw[thick] (4.8\dx,2\dx) -- (5.2\dx,2\dx) node{\rlap{$y_i$}};}
  \visible<3,4,6>{\draw[thick] (4.8\dx,4\dx) -- (5.2\dx,4\dx) node{\rlap{$f_i$}};}
  \visible<5    >{\draw[thick] (4.8\dx,3\dx) -- (5.2\dx,3\dx) node{\rlap{$f_{i+1}$}};}
  % braces
  \visible<3,6>{\draw[brace] (-2\dx, 4.2\dx) -- node[anchor=east,xshift=-0.1\dx]{\llap{\ulab{infections}}}    (-2\dx, 2.8\dx);}
  \visible<2,6>{\draw[brace] (-2\dx, 2.2\dx) -- node[anchor=east,xshift=-0.1\dx]{\llap{\ulab{symptom onset}}} (-2\dx, 0.8\dx);}
\end{tikzpicture}
}
\end{frame}
% ------------------------------------------------------------------------------
\begin{frame}{Recovering the Generation Time Distribution}
  Random variables:
  \begin{equation*}
    \left.
    \begin{aligned}
      g_i &= [f_{i+1} - f_{i}] &&\sim G(\tau) \\
      s_i &= [s_{i+1} - s_{i}] &&\sim S(\tau) \\
      h_i &= [s_{i}   - f_{i}] &&\sim H(\tau) \\
    \end{aligned}
    \visible<2->{\enspace\right\}\enspace}
    \begin{aligned}\\
      \visible<2->{s_i &= g_{i} + h_{i+1} - h_{i}}\\
      \visible<3->{    &\sim G(\tau) \conv H(\tau) \conv H(-\tau)}
    \end{aligned}
  \end{equation*}
  \pause\pause
  \invisible{\footcite{Hogg2005}}
  where $\conv$ is convolution.\pause Thus:
  \begin{equation*}
    G(\tau) = \big[S(\tau) \iconv H(\tau)\big] \iconv H(-\tau)
  \end{equation*}
\end{frame}
% ------------------------------------------------------------------------------
\begin{frame}{Practical Deconvolution}
  \begin{itemize}
    \item True deconvolution: unstable
    \item Solution: assume parametric $\hat{G}(\tau)$:
    \begin{equation*}
      \hat{G}(\tau\,|\,\theta) \sim \mathsf{Gamma}(\tau\,|\,\theta)
    \end{equation*}
    \pause
    \item Estimate $\theta$ to minimize:
    \begin{equation*}
      D_{KL} \Big( S(\tau) \,\Big|\Big|\, \hat{S}(\tau\mid\theta) \Big)
      ,\quad
      \hat{S}(\tau\,|\,\theta) = \hat{G}(\tau\,|\,\theta) \conv H(\tau) \conv H(-\tau)
    \end{equation*}
  \end{itemize}
\end{frame}
% ------------------------------------------------------------------------------
\begin{frame}{Estimating $G(\tau)$: Data}
  \heading{What is the Generation Time distribution for \covid?}
  \begin{itemize}
    \item Serial interval: \textcite{Du2020}, $N = 468$
    \begin{equation*}
      S(\tau) \sim \mathsf{Norm\,(\mu=3.96, \sigma=4.75)}
    \end{equation*}
    \item Incubation period: \textcite{Lauer2020}, $N = 181$
    \begin{equation*}
      H(\tau) \sim \mathsf{Gamma}\,(\alpha=5.81, \beta=0.95)
    \end{equation*}
  \end{itemize}
\end{frame}
% ------------------------------------------------------------------------------
\begin{frame}{\covid \& Pre-symptomatic Transmission}
  \begin{itemize}
    \item Pre-symptomatic transmission $\rightarrow$ $S(\tau)$ can be negative
    \item However: $G(\tau)$ must be non-negative
  \end{itemize}
  \pause
  \heading{Popular R package \texttt{EpiEstim}}
  \begin{itemize}
    \item Estimates $R_e(t)$ by renewal equation
    \item Input: ``serial interval'', but must be non-negative
  \end{itemize}
\end{frame}
% ------------------------------------------------------------------------------
\begin{frame}{Estimating $R_e(t)$: Data}
  \heading{What is the potential bias in $R_e(t)$ using $S(\tau)$ vs $G(\tau)$?}
  \begin{itemize}
    \item $I(t)$: \covid in GTA during March--April 2020\\[1em]
    \pause
    \item $G(\tau)$: ``recovered''\\[1em]
    \item $S(\tau)$: from literature:
    \begin{itemize}
      \item $\mathsf{Norm\,(\mu=3.96, \sigma=4.75)}$    \ttl{.35}\textcite{Du2020}\ttl{.6}Negative permitting
      \item $\mathsf{Gamma\,(\alpha=3.62, \beta=1.42)}$ \ttl{.35}\textcite{Zhang2020}\ttl{.6}Non-Negative
      \item $\mathsf{LogN\,(\mu=4.7, \sigma=2.9)}$      \ttl{.35}\textcite{Nishiura2020}\ttl{.6}Non-Negative
    \end{itemize}
  \end{itemize}
  \invisible{\footcite{Knight2020}}
\end{frame}
% ==============================================================================
\section{Results}
% ------------------------------------------------------------------------------
\begin{frame}{Recovered Generation Time Distribution}
  \begin{minipage}{0.65\linewidth}
    \includegraphics[width=\linewidth]{deconv}
  \end{minipage}\hfill%
  \begin{minipage}{0.33\linewidth}
    \pause
    $\hat{G}(\tau\,|\,\alpha,\beta)$\\
    $\sim \mathsf{Gamma}\,(\alpha=1.63,\beta=2.50)$
    \vspace{1em}
    \begin{itemize}
      \item mean: \ttl{.25} 4.08 \ttl{.5}$\approx$ 3.98
      \item SD:   \ttl{.25} 3.19 \ttl{.5}$<$ 4.75
    \end{itemize}
  \end{minipage}%
\end{frame}
% ------------------------------------------------------------------------------
\begin{frame}{Generation Time vs Serial Interval Distributions}
    \begin{minipage}{0.65\linewidth}
    \includegraphics[width=\linewidth]{distr-refs}
  \end{minipage}\hfill%
  \begin{minipage}{0.33\linewidth}
    \setlength{\tabcolsep}{3pt}
    \begin{tabular}{lcll}
      \toprule
      Ref.     & Dist. & Mean & SD \\
      \midrule
      Du       & $S(\tau)$ & 3.96 & 4.75 \\
      Zhang    & $S(\tau)$ & 5.1 & 2.7 \\
      Nishiura & $S(\tau)$ & 4.7 & 2.9 \\
      (this)   & $G(\tau)$ & 4.08 & 3.19 \\
      \bottomrule
    \end{tabular}
  \end{minipage}
\end{frame}
% ------------------------------------------------------------------------------
\begin{frame}{\covid Infections $I(t)$ in GTA}
  \begin{minipage}{0.65\linewidth}
    \includegraphics[width=\linewidth]{cases}
  \end{minipage}\hfill%
  \begin{minipage}{0.33\linewidth}
    
  \end{minipage}
\end{frame}
% ------------------------------------------------------------------------------
\begin{frame}{$R_e(t)$ using $S(\tau)$ vs $G(\tau)$}
  \begin{minipage}{0.65\linewidth}
    \includegraphics[width=\linewidth]{Re}
  \end{minipage}\hfill%
  \begin{minipage}{0.33\linewidth}
    \pause
    Compared to $G(\tau)$:
    \begin{itemize}
      \visible<3->{\item {\color{teal}Non-Negative} $S(\tau)$:\\
      \textbf{overestimate} $R_e(t)$\\
      due to: larger mean}
      \visible<4->{\item {\color{red}Negative-Permitting} $S(\tau)$:\\
      \textbf{underestimate} $R_e(t)$\\
      due to: larger variance}
    \end{itemize}
  \end{minipage}
\end{frame}
% ==============================================================================
\section{Implications}
% ------------------------------------------------------------------------------
\begin{frame}{Implications}
  \pause
  \heading{1. Characterized generation time distribution for \covid}
  \begin{itemize}
    \item Generation time\ \ =\ \ distribution of infectiousness (after exposure)
  \end{itemize}
  \pause
  \heading{2. New method to infer $G(\tau)$ from $S(\tau)$}
  \begin{itemize}
    \item Applicable to other infections
  \end{itemize}
  \pause
  \heading{3. Illustrated potential bias in $R_e(t)$ using $S(\tau)$ vs $G(\tau)$}
  \begin{itemize}
    \item Non-negative $S(\tau)$: overestimate $R_e(t)$
    \item Negative-permitting $S(\tau)$: underestimate $R_e(t)$
  \end{itemize}
\end{frame}
% ------------------------------------------------------------------------------
\begin{frame}{Limitations}
  \begin{itemize}
    \pause
    \item Assumed independent distributions $G(\tau)$, $S(\tau)$, $H(\tau)$
    \pause
    \item Parametric approximations\\
          However: useful for meta-analyses
    \pause
    \item No uncertainty analysis / confidence intervals
  \end{itemize}
\end{frame}
% ------------------------------------------------------------------------------
\begin{frame}{Thanks}
  \centering\footnotesize
  \heading{\large Funding}\\
  \begin{tabularx}{\linewidth}{CCC}
    \includegraphics[height=0.10\textheight]{nserc} &
    \includegraphics[height=0.10\textheight]{on} &
    \includegraphics[height=0.08\textheight]{smh} \\
    Canada Graduate Scholarship--D & Ontario Early Researcher Award & 2020 \covid Research Award
  \end{tabularx}
  \vfill
  \heading{\large Support}\\
  \begin{tabularx}{.5\linewidth}{CC}
    \includegraphics[height=0.12\textheight]{uoft} &
    \includegraphics[height=0.10\textheight]{map}
  \end{tabularx}
  \vfill
  Kristy Yiu, Linwei Wang, Huiting Ma, David Landsman, \textit{Unity Health Toronto}\\
  David Champredon, \textit{University of Western Ontario}\\
  Elena Aruffo, Jude Dzevela Kong, \textit{CAIMS 2021 Mini Symposium Organizers}
\end{frame}
% ------------------------------------------------------------------------------
\begin{frame}{References}
  \printbibliography[heading=none]
\end{frame}
% ------------------------------------------------------------------------------
\begin{frame}{Preliminary Uncertainty Analysis}
  \begin{minipage}{0.65\linewidth}
    \centering
    \includegraphics[height=.85\textheight]{uncertainty}
  \end{minipage}\hfill%
  \begin{minipage}{0.33\linewidth}
    Likelihood of $\alpha, \beta$ for $\hat{G}(\tau)$\\
    from ``deconvolution''
  \end{minipage}
\end{frame}
% ==============================================================================
\end{document}